
\documentclass{article}
\usepackage{colortbl}
\usepackage{makecell}
\usepackage{multirow}
\usepackage{supertabular}

\begin{document}

\newcounter{utterance}

\centering \large Interaction Transcript for game `wordle', experiment `high\_frequency\_words\_no\_clue\_no\_critic', episode 3 with Qwen3.6{-}35B{-}A3B{-}FP8{-}without{-}reasoning.
\vspace{24pt}

{ \footnotesize  \setcounter{utterance}{1}
\setlength{\tabcolsep}{0pt}
\begin{supertabular}{c@{$\;$}|p{.15\linewidth}@{}p{.15\linewidth}p{.15\linewidth}p{.15\linewidth}p{.15\linewidth}p{.15\linewidth}}
    \# & \multicolumn{2}{c}{Player} && \multicolumn{2}{c}{Game Master} \\
    \hline

    \theutterance \stepcounter{utterance}  
    & & & \multicolumn{4}{p{0.6\linewidth}}{
        \cellcolor[rgb]{0.9,0.9,0.9}{
            \makecell[{{p{\linewidth}}}]{
                \texttt{\tiny{[P1$\langle$GM]}}
                \texttt{Du bist ein Sprachzauberer, der gerne Wörter anhand der gegebenen Regeln errät.} \\
\\ 
\texttt{Willkommen bei Wordle! Du hast sechs Versuche, um das Zielwort zu erraten, ein gültiges deutsches Wort aus fünf Kleinbuchstaben. Bitte verwende die Tags "erklärung:" und "vermutung:", um für jeden Tipp eine knappe Erklärung zu geben.} \\
\\ 
\texttt{Wenn dein Tipp zum Beispiel "Firma" ist, sollte deine Antwort lauten} \\
\texttt{erklärung: dies ist ein gängiges deutsches Wort mit fünf Buchstaben, und ich beginne meinen Tipp mit diesem Wort.} \\
\texttt{vermutung: Firma} \\
\\ 
\texttt{Nach jedem Tipp wird deine Antwort überprüft, und du erhältst Feedback, das anzeigt, welche Buchstaben korrekt sind (grün), welche Buchstaben korrekt, aber an der falschen Position sind (gelb), und welche Buchstaben falsch sind (rot). Dieses Feedback kann hilfreich sein, um zu bestimmen, welche Buchstaben du in deinem nächsten Tipp einbeziehen oder ausschließen solltest.} \\
\\ 
\texttt{Zum Beispiel könnte das Feedback für "Firma" so aussehen:} \\
\texttt{rückmeldung\_zur\_vermutung: f<gelb> i<gelb> r<grün> m<gelb> a<rot>} \\
\\ 
\texttt{Die Erklärung sollte Details dazu enthalten, wie das rückmeldung\_zur\_vermutung verwendet wird, um zu einem neuen Tipp zu gelangen.} \\
\\ 
\texttt{Erzeuge keinen anderen Text.} \\
\\ 
\texttt{Beginnen wir mit deinem ersten Tipp.} \\
            }
        }
    }
     \\ \\

    \theutterance \stepcounter{utterance}  
    & \multicolumn{4}{p{0.6\linewidth}}{
        \cellcolor[rgb]{0.561,0.737,0.561}{
            \makecell[{{p{\linewidth}}}]{
                \texttt{\tiny{[P1$\rangle$GM]}}
                \texttt{erklärung: Ich beginne mit einem sehr häufigen deutschen Fünf{-}Buchstaben{-}Wort, das die Vokale A, E und I enthält, um schnell viel Information über die Vokale und einige Konsonanten zu erhalten.} \\
\texttt{vermutung: reise} \\
            }
        }
    }
    & & \\ \\

    \theutterance \stepcounter{utterance}  
    & & & \multicolumn{2}{p{0.3\linewidth}}{
        \cellcolor[rgb]{0.9,0.9,0.9}{
            \makecell[{{p{\linewidth}}}]{
                \texttt{\tiny{[GM$|$GM]}}
                \texttt{attempts = 1} \\
\texttt{target\_word = april} \\
\texttt{guess = reise} \\
\texttt{guess\_feedback = r<yellow> e<red> i<yellow> s<red> e<red>} \\
            }
        }
    }
    & & \\ \\

    \theutterance \stepcounter{utterance}  
    & & & \multicolumn{4}{p{0.6\linewidth}}{
        \cellcolor[rgb]{0.9,0.9,0.9}{
            \makecell[{{p{\linewidth}}}]{
                \texttt{\tiny{[P1$\langle$GM]}}
                \texttt{rückmeldung\_zur\_vermutung: r<gelb> e<rot> i<gelb> s<rot> e<rot>} \\
\\ 
\texttt{Gib deine Antwort nur in diesem Format an.} \\
\texttt{erklärung: details} \\
\texttt{vermutung: wort} \\
            }
        }
    }
     \\ \\

    \theutterance \stepcounter{utterance}  
    & \multicolumn{4}{p{0.6\linewidth}}{
        \cellcolor[rgb]{0.561,0.737,0.561}{
            \makecell[{{p{\linewidth}}}]{
                \texttt{\tiny{[P1$\rangle$GM]}}
                \texttt{erklärung: Das Feedback zeigt, dass 'r' und 'i' im Wort vorkommen, aber nicht an den Positionen 1 bzw. 3. 'e' und 's' sind ausgeschlossen. Ich versuche nun, 'r' und 'i' an anderen Positionen unterzubringen und füge neue Konsonanten hinzu.} \\
\texttt{vermutung: trier} \\
            }
        }
    }
    & & \\ \\

    \theutterance \stepcounter{utterance}  
    & & & \multicolumn{2}{p{0.3\linewidth}}{
        \cellcolor[rgb]{0.9,0.9,0.9}{
            \makecell[{{p{\linewidth}}}]{
                \texttt{\tiny{[GM$|$GM]}}
                \texttt{attempts = 2} \\
\texttt{target\_word = april} \\
\texttt{guess = trier} \\
\texttt{guess\_feedback = t<red> r<yellow> i<yellow> e<red> r<red>} \\
            }
        }
    }
    & & \\ \\

    \theutterance \stepcounter{utterance}  
    & & & \multicolumn{4}{p{0.6\linewidth}}{
        \cellcolor[rgb]{0.9,0.9,0.9}{
            \makecell[{{p{\linewidth}}}]{
                \texttt{\tiny{[P1$\langle$GM]}}
                \texttt{rückmeldung\_zur\_vermutung: t<rot> r<gelb> i<gelb> e<rot> r<rot>} \\
\\ 
\texttt{Gib deine Antwort nur in diesem Format an.} \\
\texttt{erklärung: details} \\
\texttt{vermutung: wort} \\
            }
        }
    }
     \\ \\

    \theutterance \stepcounter{utterance}  
    & \multicolumn{4}{p{0.6\linewidth}}{
        \cellcolor[rgb]{0.561,0.737,0.561}{
            \makecell[{{p{\linewidth}}}]{
                \texttt{\tiny{[P1$\rangle$GM]}}
                \texttt{erklärung: 'r' ist an Position 2 korrekt (grün), 'i' ist im Wort, aber nicht an Position 3 (gelb). 't' und 'e' sind ausgeschlossen. Ich platziere 'i' an Position 1 oder 4 und versuche neue Konsonanten.} \\
\texttt{vermutung: irren} \\
            }
        }
    }
    & & \\ \\

    \theutterance \stepcounter{utterance}  
    & & & \multicolumn{2}{p{0.3\linewidth}}{
        \cellcolor[rgb]{0.9,0.9,0.9}{
            \makecell[{{p{\linewidth}}}]{
                \texttt{\tiny{[GM$|$GM]}}
                \texttt{attempts = 3} \\
\texttt{target\_word = april} \\
\texttt{guess = irren} \\
\texttt{guess\_feedback = i<yellow> r<red> r<green> e<red> n<red>} \\
            }
        }
    }
    & & \\ \\

    \theutterance \stepcounter{utterance}  
    & & & \multicolumn{4}{p{0.6\linewidth}}{
        \cellcolor[rgb]{0.9,0.9,0.9}{
            \makecell[{{p{\linewidth}}}]{
                \texttt{\tiny{[P1$\langle$GM]}}
                \texttt{rückmeldung\_zur\_vermutung: i<gelb> r<rot> r<grün> e<rot> n<rot>} \\
\\ 
\texttt{Gib deine Antwort nur in diesem Format an.} \\
\texttt{erklärung: details} \\
\texttt{vermutung: wort} \\
            }
        }
    }
     \\ \\

    \theutterance \stepcounter{utterance}  
    & \multicolumn{4}{p{0.6\linewidth}}{
        \cellcolor[rgb]{0.561,0.737,0.561}{
            \makecell[{{p{\linewidth}}}]{
                \texttt{\tiny{[P1$\rangle$GM]}}
                \texttt{erklärung: 'r' ist an Position 3 korrekt (grün). 'i' ist im Wort, aber nicht an Position 1 (gelb). 'r' an Position 2 ist falsch, also ist das zweite 'r' in 'irren' nicht korrekt platziert oder das Wort hat nur ein 'r' an Pos 3. Da 'r' an Pos 3 grün ist, muss das Wort die Struktur \_ \_ r \_ \_ haben. 'i' ist enthalten, aber nicht an Pos 1. 'e' und 'n' sind ausgeschlossen. Ich versuche 'i' an Pos 2 oder 4 und neue Buchstaben.} \\
\texttt{vermutung: irisch} \\
            }
        }
    }
    & & \\ \\

    \theutterance \stepcounter{utterance}  
    & & & \multicolumn{2}{p{0.3\linewidth}}{
        \cellcolor[rgb]{0.9,0.9,0.9}{
            \makecell[{{p{\linewidth}}}]{
                \texttt{\tiny{[GM$|$GM]}}
                \texttt{Error: The length of the guessed word is not 5.} \\
            }
        }
    }
    & & \\ \\

    \theutterance \stepcounter{utterance}  
    & & & \multicolumn{2}{p{0.3\linewidth}}{
        \cellcolor[rgb]{0.9,0.9,0.9}{
            \makecell[{{p{\linewidth}}}]{
                \texttt{\tiny{[GM$|$GM]}}
                \texttt{game\_result = ABORT} \\
            }
        }
    }
    & & \\ \\

\end{supertabular}
}

\end{document}
